\subsection{Choose 16 kHz with a measured cost comparison}
The target is the complete supplied multitone, including its weak highest
component. The reference sampling condition and replica geometry are
 \citep{shannon1949noise}
\begin{equation}
 T_s=1/f_s,\quad x_s[m]=x(mT_s),\quad
 f_s>2f_B=12000\,\mathrm{Hz},\quad f_N=f_s/2.
 \label{eq:rate-condition}
\end{equation}
\begin{equation}
 g=f_s/2-f_B=\SamplingGuard\,\mathrm{Hz},\qquad
 d=f_s-2f_B=2g=\ReplicaGap\,\mathrm{Hz}.
 \label{eq:rate-guard}
\end{equation}
The signal's Nyquist rate is 12 kHz; the selected sampler's Nyquist
frequency is 8 kHz. The 2 kHz guard is a sampling margin, not a measured
filter specification. Following a predict--sample--compare method,
I take actual reference values, not newly generated ideal tones:
\begin{equation}
 x_M[m]=x_{\mathrm{ref}}[Mm],\quad
 f_s=F_0/M,\quad N=L/M,\quad
 t_m=\frac m{f_s}=\frac{Mm}{F_0}.
 \label{eq:integer-sampling}
\end{equation}
$ M=4,3,2,1 $ tests 12, 16, 24, 48 kHz. All records span 2 s and have 0.5 Hz native
spacing. The implementation is \texttt{reference[::factor].copy()}.
The phase-aware folding identity is Eq.~\eqref{eq:decimation-folding};
an independent full-complex-FFT comparison of every bin at five measured
rates differs by at most $\FoldingError $ NPY units.

\begin{table}[H]
 \centering
 \caption{Same 2 s record, common periodic recovery and all 96000 evaluation positions. Raw payload compares equal bits/sample against 16 kHz; 12 kHz is phase-sensitive critical sampling.}
 \label{tab:rates}
 \begin{tabular}{rrrrrr}
\toprule
$f_s$ (kHz) & $M$ & $f_N$ (kHz) & Samples & Raw ratio & MSE \\
\midrule
12 & 4 & 6 & 24,000 & 0.75 & $0.005$ \\
16 & 3 & 8 & 32,000 & 1 & $1.75\times10^{-25}$ \\
24 & 2 & 12 & 48,000 & 1.5 & $1.27\times10^{-25}$ \\
48 & 1 & 24 & 96,000 & 3 & $1.72\times10^{-31}$ \\
\bottomrule
\end{tabular}

\end{table}

Reconstruction uses only low-rate samples and output length, with the
periodic Fourier interpolator of Eq.~\eqref{eq:periodic-interpolation}.
Its $ L/N $ scale and split old-Nyquist term are derived in
Appendix~\ref{app:sampling}. Coherent 2 s periodicity supports this boundary
model. DC/phased-cosine/Nyquist calibrations and SciPy resampling check it
 \citep{scipy1153resample}. Every metric uses all 96000 positions
(Eq.~\eqref{eq:sampling-errors}); no trimming, fitted shift/gain or
reference-assisted recovery is used. At 16 kHz, MSE/NMSE are
$\SamplingMse/\SamplingNmse $, peak error $\SamplingMaxError $.

The 12 kHz boundary loses the 6 kHz sine through its measured phase near
$-\pi/2 $ (Eq.~\eqref{eq:critical-phase}): endpoint amplitude
$\CriticalToneAmplitude $, MSE $\CriticalMse $, agreeing with the missing
tone power 0.005. Higher 24/48 kHz rates have numerical-residual errors too,
but require more payload. Since $48000/M>12000 $ implies $ M<4 $,
16 kHz is the lowest sufficient tested integer-decimation choice;
it saves one third/two thirds of 24/48 kHz raw bits at fixed $ B $
(Eq.~\eqref{eq:sampling-payload}). The theoretically valid 14.4 kHz
alternative would save 10 \% relative to 16 kHz with 1.2 kHz guard, but needs
noninteger resampling, which I have not measured. I accept 11.11 \% extra
payload relative to it for direct indexing and a larger guard.
This is not a global optimum under later noise.

No appreciable source components lie beyond 6 kHz, so no extra prefilter
is inserted in this clean index-selection study. An arbitrary broadband
input would require anti-alias filtering. The specified noise is added
after downsampling and cannot alias through that earlier operation.
Twelve calibration/grid/phase/error checks pass; R1 is deterministic.
